\section{The 209-state and 196-state published designs}

The comparison uses the PA209 commit
\texttt{69d98241a721ead20b4f0b17df9062950b40fbfe} and the PA196 commit
\texttt{407cdc356f6b586635e71ac92be394255f4bfea4} of
\url{https://github.com/jamesyc/BusyBeaver}. The latter is two commits after
the former; commit \texttt{64ff1b3} replaces the proof package and
\texttt{407cdc3} polishes its paper.

\begin{center}
\begin{tabular}{lrr}
\hline
Quantity & PA209 & PA196\\\hline
Working states &209&196\\
Defined binary transition rows&418&392\\
Source instruction sites&3948&3812\\
Distinct internal source-tree nodes&173&159\\
Maximum instruction-path length&33&38\\
Startup steps to the candidate-2 frame&58210&58154\\
AX macro tokens under the common tokenization&210&201\\
AX internal grammar nodes under that tokenization&83&68\\\hline
\end{tabular}
\end{center}
The last two rows are independently counted from the published Lean source
trees. A token is one of push, push-and-copy of one of the three inputs,
the Cantor constructor, increment, the parameter readers, or the final
guard/control macros. The distinct-tree count includes every reachable
internal node, not just nodes with names in the generated source.

PA209 uses sparse labels $0,\ldots,213$ excluding
$50,122,150,157,205$, with halt label $215$. PA196 uses consecutive
working labels $0,\ldots,195$ and halt label $196$. State renaming alone
therefore obscures the useful structural comparison.

\subsection{Logical and coding changes}

Writing $p$ for Cantor pairing, PA209 uses $4p(a,b)+r$ for all four
constructors, with residues $0,1,2,3$ for implication, membership,
universal quantification, and equality. PA196 retains the implication,
membership, and equality constructors but uses
$U(a,B)=4p(B,a)+2$. This transposition makes repeated quantifier
construction share more of the register-program grammar.

The decisive logical change replaces the seventh primitive-equality
head by guarded equality vacuity
$E(a,b)\to U(c,E(a,b))$. PA196 adds that clause to the separation gate,
under its existing freshness guard $c\ne a,b$. Reverse membership
rewriting is recovered using pointed separation; the interpretation of
arithmetic then uses equality defined by matching both incoming and
outgoing membership neighborhoods. It does not need primitive equality
symmetry and transitivity as independent axioms.

The remaining six heads are reordered, two pairs of implication
antecedents are exchanged, and the equality witness is mirrored.
These changes improve repeated-subtree sharing while preserving the
required interpreted equality facts. The 196-state machine is therefore
not simply the old table with thirteen duplicate rows removed.

\subsection{Physical sharing}

PA196 has 36 engine states, 13 leaf-routing states, and 147 decision
states for its 159 distinct source-tree nodes. Six pairs of AX grammar
nodes share controls using the bit immediately preceding the current
path bit. For a conflicting branch, the shared control steps left,
reads this saved choice, and uses existing controls to step back right
into the correct continuation. These six reductions do not introduce
additional helper states. The other existing control folds remain in
the unchanged outer framework. Two obsolete outside-body distinction
states of PA209 are removed.

An independent compiler in \texttt{compile\_ax.py} reconstructs the
unfolded 202-state PA196 decoder and reproduces the supplied 196-state
transition table exactly after the published six AX folds. Thus its
state-count model is checked against the literal machine.

\subsection{What the published formal endpoints assert}

Both packages explicitly prove an internal implication
\[
  \mathrm{PA}\vdash\neg H\to\operatorname{Con}(\mathrm{PA})
\]
and an external consistency-to-nonhalting theorem, together with
two-sided independence under the stated metatheoretic assumptions.
Their README files and theorem tables do not advertise a checked
internal converse
$\mathrm{PA}\vdash\operatorname{Con}(\mathrm{PA})\to\neg H$.
That converse requires an arithmetized syntactic soundness compiler;
it must not be inferred merely from the external theorem.

\section{Conservative changes to the command interpreter}

Two independently checked program changes produce alternative
196-state tables. They can be combined with further decoder
optimizations, but by themselves do not satisfy the strict target.

First, a successful MP command need not push a zero onto its stack.
Changing native row $(164,1)$ from target $98$ to target $114$
(dense row $(115,1)$ from target $86$ to target $100$) removes the
19-site push block. The new source has 3793 sites. Resolving every
remaining source instruction gives exactly the old control-flow graph
with that block deleted and addresses shifted. All concrete primitive
instruction contracts checked at clean frames, including both outcomes
of each decrement, pass. At the command level MP success becomes
$(b,w,b,0)$ rather than $(b,p(0,w),b,0)$. A refutation compiler can
insert zero below the two premises using the old discard-and-push
words before invoking the otherwise unchanged MP compilation.

Second, a successful MP command may destructively move its consequent
from register $D$ into visible register $H$. Native row $(164,1)$ then
targets $130$ (dense target $108$), and the source replaces the
push-and-copy subtree by the already present three-instruction
$D\to H$ loop. The resulting source has 3789 sites. All 4928 concrete
primitive instruction cases pass. MP success is $(0,w,0,b)$.
With a zero inserted below the premises, the word consisting of two
Shifts, MP$(b)$, and one Shift returns $(b,w,0,0)$, or halts when
$b=1$. This supplies a particularly simple adaptation of the
refutation compiler.

These concrete runs are falsification checks, not a substitute for a
uniform proof over all natural register values. The supplied source
and row changes make the required uniform proof local: the standard
unary engine is unchanged, and the replacement loop has the elementary
invariant that $D+H$ is constant until $D=0$.
For both changed tables, a separate exact finite computation confirms
the complete original startup configuration after 58154 steps:
candidate two, gap 48, and the physical file
$[0,0,0,0,0,0,2,0,0]$ followed by 26 ones.

\section{A checked family of physical register permutations}

The logical registers are physically ordered by the permutation
$[6,2,0,4,3,7,8,5,1]$. Retargeting the eleven non-engine rows which enter
the fixed nine-state register-skip chain permutes these physical
locations. Keeping the candidate counter at physical position six gives
28 elementary transpositions of the other eight positions. All 28
tables pass all 4958 primitive instruction cases in the clean-frame
test profile, for a total of 138824 checked runs. None yielded a further
state fold in the common-row finite-window search at radius four.
This is a negative result for the specified finite search, not a
lower bound on state complexity.
The same 28 transpositions were also combined with the destructive
MP variant and screened at radius four; no further fold was found.
Those combined-table screens are not additional uniform simulation
proofs.

Changing the physical order can also change bootstrap behavior.
For the transposition of physical positions zero and one, an exact
blank-tape run reaches candidate two after 43248 steps with gap 41
and 28 registers. That gap is below the published uniform bound.
After 74404 steps it reaches a fully checked clean candidate-three
frame with gap 53 and 41 registers; every other logical register is
zero and the 32 spare registers are one. This later frame meets the
published gap bound. One should therefore record the actual new
startup frame, rather than expect the old 58154-step trace to survive
a register permutation unchanged. The source scheduler's completeness
argument can start at candidate three, since the compiled universal
halting word has candidate index at least three.
